\section*{第 \ref{ch:b1:cell-unit-of-life} 章——细胞：生命的基本单位}

\begin{solution}{exo:b1:cell-unit-of-life:1}
所有\omterm{def:b1:organism-environment:organism}{生物体}都由\omterm{def:b1:cell-unit-of-life:cell}{细胞}构成（在显微镜下检查的每一种\omterm{def:b1:body-plans-tissues:tissue}{组织}，
无论动物还是植物，都可见\omterm{def:b1:cell-unit-of-life:cell}{细胞}）；每个\omterm{def:b1:cell-unit-of-life:cell}{细胞}都来自另一个\omterm{def:b1:cell-unit-of-life:cell}{细胞}
（在胚胎和藻类中观察到了分裂；巴斯德的烧瓶排除了自发生成）；
遗传物质在分裂时传递
（染色体可通过有丝分裂被追踪）。
\end{solution}

\begin{solution}{exo:b1:cell-unit-of-life:2}
\omterm{def:b1:cell-unit-of-life:prokeuk}{细胞核}（无 / 有），有膜包被的\omterm{def:b1:cell-unit-of-life:prokeuk}{细胞器}（无 / 有），
染色体（一条环状 / 多条线性），大小
（\qtyrange{0.5}{5}{\micro m} / \qtyrange{10}{100}{\micro m}），
\omterm{def:b1:cell-unit-of-life:prokeuk}{细胞壁}由肽聚糖构成（细菌）/ 由纤维素、几丁质构成或无，
\omterm{def:b1:cell-unit-of-life:cell}{细胞}骨架（不发达 / 发达）。
\end{solution}

\begin{solution}{exo:b1:cell-unit-of-life:3}
$\qty{1}{mm}/\qty{1}{\micro m} = 1000$ 个细菌；$1000/20 = 50$ 个动物
\omterm{def:b1:cell-unit-of-life:cell}{细胞}。
\end{solution}

\begin{solution}{exo:b1:cell-unit-of-life:4}
\omterm{def:b1:cell-unit-of-life:resolution}{分辨率}取决于波长和数值孔径，而非\omterm{def:b1:cell-unit-of-life:resolution}{放大}倍数；
超过约 $1000\times$ 之后，图像只是\omterm{def:b1:cell-unit-of-life:resolution}{放大}了衍射极限
（\qty{0.2}{\micro m}）造成的模糊，
并未\omterm{def:b1:cell-unit-of-life:resolution}{分辨}出任何新的细节。
\end{solution}

\begin{solution}{exo:b1:cell-unit-of-life:5}
$0.61\times 550/0.65 = \qty{516}{nm}$；若用蓝光则为 \qty{422}{nm}。两
个相距 \qty{500}{nm} 的细菌：勉强可以在蓝光下\omterm{def:b1:cell-unit-of-life:resolution}{分辨}；核糖体
相距 \qty{30}{nm}：不能\omterm{def:b1:cell-unit-of-life:resolution}{分辨}，相差约十五倍。
\end{solution}

\begin{solution}{exo:b1:cell-unit-of-life:6}
紫色：较厚的肽聚糖壁保留了结晶紫–碘复合物，
无外膜（革兰氏阳性）。粉红色：较薄的肽聚糖层
位于外膜之下，染料被洗脱，经复染（革兰氏阴性）。
\emph{E. coli} 是革兰氏阴性杆菌：粉红色的杆菌。
\end{solution}

\begin{solution}{exo:b1:cell-unit-of-life:7}
线粒体：$v = \frac{2}{9}\times(5\times 10^{-7})^2\times 100\times
9.81\times 10^4/10^{-3} = \qty{5.4e-4}{m/s}$；\qty{5}{cm} 沉降
需 \qty{92}{s}。\omterm{def:b1:cell-unit-of-life:prokeuk}{细胞核}在 $1000\,g$ 下：$v = \frac{2}{9}\times(4\times
10^{-6})^2\times 100\times 9810/10^{-3} = \qty{3.5e-3}{m/s}$; \qty{5}{cm}
需 \qty{14}{s}。
\end{solution}

\begin{solution}{exo:b1:cell-unit-of-life:8}
琥珀酸脱氢酶存在于线粒体中：\qty{80}{\%} 的
线粒体位于沉淀 2 中；沉淀 1 中的 \qty{15}{\%} 是
被困在未破碎\omterm{def:b1:cell-unit-of-life:cell}{细胞}和碎片中、或随\omterm{def:b1:cell-unit-of-life:prokeuk}{细胞核}一同沉降的线粒体；
上清液中的 \qty{5}{\%} 来自\omterm{def:b1:cell-unit-of-life:fractionation}{匀浆}时被破坏的线粒体。乳酸脱氢酶
存在于\omterm{def:b1:cell-unit-of-life:cell}{细胞}溶胶中，是可溶的，
因而留在最终上清液中。
\end{solution}

\begin{solution}{exo:b1:cell-unit-of-life:9}
(a) 活\omterm{def:b1:cell-unit-of-life:cell}{细胞}的相差\omterm{prop:b1:cell-unit-of-life:microscopes}{光学显微镜}观察（该过程
需数分钟，并且需要活\omterm{def:b1:cell-unit-of-life:cell}{细胞}）。(b) 染色切片的\omterm{prop:b1:cell-unit-of-life:microscopes}{光学显微镜}观察，
或为确证而用透射\omterm{prop:b1:cell-unit-of-life:microscopes}{电子显微镜}（线粒体大小恰处于\omterm{prop:b1:cell-unit-of-life:microscopes}{光学显微镜}\omterm{def:b1:cell-unit-of-life:resolution}{分辨}极限）。
(c) 透射\omterm{prop:b1:cell-unit-of-life:microscopes}{电子显微镜}（孔的大小为 \qty{100}{nm}）。(d) 扫描\omterm{prop:b1:cell-unit-of-life:microscopes}{电子显微镜}（表面起伏）。
\end{solution}

\begin{solution}{exo:b1:cell-unit-of-life:10}
清澈的烧瓶排除了这样一种情形：空气本身或煮沸的肉汤本身
会产生生命；空气可以经敞开的瓶颈自由进入。浑浊的烧瓶
与尘埃接触后表明，瓶颈所截留的物质——来自空气的颗粒——
才是生物的来源。
两者共同证明，微生物来自被尘埃携带的微生物。
密封后煮沸的烧瓶若敞口放置，会遭到这样的反驳：密封
排除了“生命力”所需的空气；天鹅颈瓶允许空气进入，
只排除尘埃。
\end{solution}

\begin{solution}{exo:b1:cell-unit-of-life:11}
$10^5 = 2^{n}$ 得 $n = 16.6$ 次倍增，即 \qty{415}{min}，约
\qty{7}{h}. $10^8\times 10^3\,\mathrm{mL}\times\qty{1}{pg} = \qty{0.1}{g}$
每升。生长停止，是因为氧气（培养基只能供应
每升几毫克）、某种无机物，或酸和废物的积累
成为限制因素，而非葡萄糖。
\end{solution}

\begin{solution}{exo:b1:cell-unit-of-life:12}
每一维都大十倍的\omterm{def:b1:cell-unit-of-life:cell}{细胞}，单位体积的膜面积只有原来的十分之一；
由于膜承载着能量转换和运输的机器，大\omterm{def:b1:cell-unit-of-life:cell}{细胞}将会“饥饿”。真核生物
通过将膜向内折叠来增加膜面积，
（内质网、线粒体内膜、类囊体）；
同时也因此形成了封闭的区室，
使互不相容的过程（酸性消化、氧化化学反应、DNA 储存）
得以在各自的条件下彼此分开进行。
就膜而言，这一比喻是贴切的，但并不完整：
真核生物还获得了\omterm{def:b1:cell-unit-of-life:cell}{细胞}骨架和\omterm{def:b1:cell-unit-of-life:prokeuk}{细胞核}，
其线粒体则是被吞入的细菌的后代，而非折叠形成的。
\end{solution}

\begin{solution}{pb:b1:cell-unit-of-life:1}
\textbf{1.} 沉淀 1：\omterm{def:b1:cell-unit-of-life:prokeuk}{细胞核}（以及未破碎的\omterm{def:b1:cell-unit-of-life:cell}{细胞}、碎片）。沉淀 2：
线粒体（含溶酶体和过氧化物酶体）。沉淀 3：微粒体
（内质网碎片）和核糖体。
\textbf{2.} 低温减缓酶的活性（包括破碎后释放的蛋白酶和
磷脂酶），并保存\omterm{def:b1:cell-unit-of-life:prokeuk}{细胞器}；
采用等渗溶液，使\omterm{def:b1:cell-unit-of-life:prokeuk}{细胞器}既不膨胀破裂，也不皱缩。
\textbf{3.} 线粒体（内膜）；\omterm{def:b1:cell-unit-of-life:cell}{细胞}溶胶。
\textbf{4.} $500 + 200 + 300 + 950 = \qty{1950}{mg}$，\qty{97.5}{\%}：
损失发生在离心管壁上以及转移过程中。
\textbf{5.} SDH：$15 + 60 + 5 + 4 = 84$，\qty{84}{\%}（部分活性
因线粒体受损而丢失）。LDH：$20 + 12 + 8 + 350 = 390$，
\qty{97.5}{\%}。
\textbf{6.} $2000/10 = \qty{200}{mg/g}$：\qty{20}{\%} 为鲜重（其余大部分是水）。
\textbf{7.} $v = \frac{2}{9}\times(5\times 10^{-7})^2\times 100\times
\num{9.81e4}/10^{-3} = \qty{5.4e-4}{m/s}$；\qty{5}{cm} 需 \qty{92}{s}：
\qty{20}{min} 绰绰有余。
\textbf{8.} 在 $1000\,g$ 下，$v = \qty{5.4e-5}{m/s}$；经过 \qty{600}{s} 后，
下沉 \qty{3.3}{cm}。若线粒体在 \qty{5}{cm} 内均匀分布，则距底部
\qty{3.3}{cm} 以内的线粒体会沉降到底部：占 \qty{65}{\%}，
前提是不受其他因素干扰。
\textbf{9.} 实际上只有 \qty{15}{\%} 形成沉淀：真实沉淀并不会
截留所有到达它的物质（紧密堆积的\omterm{def:b1:cell-unit-of-life:prokeuk}{细胞核}被温和地重新悬浮，
疏松的上层则留在上清液中），
且该公式高估了非球形、含水\omterm{def:b1:cell-unit-of-life:prokeuk}{细胞器}的沉降速度。
未破碎\omterm{def:b1:cell-unit-of-life:cell}{细胞}内的以及黏附在\omterm{def:b1:cell-unit-of-life:prokeuk}{细胞核}上的线粒体
也会进入沉淀 1。
\textbf{10.} $v = \frac{2}{9}\times(4\times 10^{-6})^2\times 100\times
9810/10^{-3} = \qty{3.5e-3}{m/s}$; \qty{14}{s}.
\textbf{11.} $v = \frac{2}{9}\times(1.2\times 10^{-8})^2\times 600\times
\num{9.81e5}/10^{-3} = \qty{1.9e-5}{m/s}$；\qty{5}{cm} 需 \qty{2650}{s}，
\qty{44}{min}。在 $\num{10000}\,g$ 下，速度为 \qty{1.9e-6}{m/s}：
经过 \qty{20}{min}，下沉 \qty{2.3}{mm}，因此核糖体留在上清液 2 中。
\textbf{12.} 速度与 $r^2$ 成正比：在较小的力作用下，只有最大的颗粒
会明显移动，因此每一步只沉降一类颗粒，留下较小的颗粒。
时间过短，该类颗粒沉降不完全；时间过长，下一类
颗粒便会污染沉淀。
\textbf{13.} \omterm{def:b1:cell-unit-of-life:fractionation}{匀浆} $100/2000 = \qty{0.05}{U/mg}$；沉淀 1
$15/500 = 0.03$；沉淀 2 $60/200 = 0.30$；沉淀 3 $5/300 = 0.017$；
上清液 $4/950 = 0.004$。
\textbf{14.} 纯化倍数 $0.30/0.05 = 6$；产率 $60/100 =
\qty{60}{\%}$。
\textbf{15.} $0.017/0.05 = 0.33$：沉淀 3 中线粒体含量低；其中少量活性来自\omterm{def:b1:cell-unit-of-life:fractionation}{匀浆}过程中产生的线粒体膜碎片，它们与微粒体一同沉降。
\textbf{16.} LDH：\omterm{def:b1:cell-unit-of-life:fractionation}{匀浆} $400/2000 = \qty{0.2}{U/mg}$；上清液
$350/950 = \qty{0.37}{U/mg}$；纯化倍数 $1.8$；产率 \qty{87.5}{\%}。
\textbf{17.} \omterm{def:b1:cell-unit-of-life:cell}{细胞}溶胶含有\omterm{def:b1:cell-unit-of-life:cell}{细胞}一半的蛋白质，因此即使是纯的
\omterm{def:b1:cell-unit-of-life:cell}{细胞}溶胶，比活性也只能提高一倍；线粒体的蛋白质含量仅为六分之一
或更少，因此分离线粒体可使其酶富集六倍或更多。
对于可溶性酶，产率较高，因为它不会
随受损的\omterm{def:b1:cell-unit-of-life:prokeuk}{细胞器}一同丢失。
\textbf{18.} 不是：沉淀 2 含有 \qty{200}{mg} 蛋白质，并在其颗粒
之间截留了上清液；\qty{12}{U} 占总量的 \qty{3}{\%}，
与截留的\omterm{def:b1:cell-unit-of-life:cell}{细胞}溶胶相符。检验方法：洗涤沉淀
（重新悬浮后再次离心）：被截留的酶会被洗去，而结合的酶则保留。
\textbf{19.} 若沉淀 2 中线粒体占 \qty{60}{\%}，则纯线粒体
的比活性为 $0.30/0.6 = \qty{0.5}{U/mg}$；肝脏的 \qty{100}{U} 则
对应 \qty{200}{mg} 线粒体蛋白质，
占 \qty{2000}{mg} 的 \qty{10}{\%}。
\textbf{20.} 对经固定、包埋、超薄切片并经重金属染色的沉淀进行透射\omterm{prop:b1:cell-unit-of-life:microscopes}{电子显微镜}观察：
可见具有双层膜和嵴的豆状轮廓，
其间夹杂着一些溶酶体（致密小体）和过氧化物酶体。
\textbf{21.} 在相差显微镜下观察一滴重新悬浮的沉淀，在游离的\omterm{def:b1:cell-unit-of-life:prokeuk}{细胞核}中计数完整\omterm{def:b1:cell-unit-of-life:cell}{细胞}；若完整\omterm{def:b1:cell-unit-of-life:cell}{细胞}较多，则应更充分地\omterm{def:b1:cell-unit-of-life:fractionation}{匀浆}（增加研磨次数，使用更紧密的研杵）并重新离心，因为被截留的\omterm{def:b1:cell-unit-of-life:prokeuk}{细胞器}会降低各项产率。
\textbf{22.} \omterm{def:b1:cell-unit-of-life:fractionation}{差速离心}按大小分离（速度
$\propto r^2$），而三类\omterm{def:b1:cell-unit-of-life:prokeuk}{细胞器}的大小相近；在
密度梯度中，颗粒会停在与其自身密度相同的位置，与大小无关，
而三类\omterm{def:b1:cell-unit-of-life:prokeuk}{细胞器}的密度则各不相同。
\textbf{23.} 完整的溶酶体将酶保留在膜内，底物无法穿过该膜（潜伏性）；
去污剂破坏了膜并使其暴露。
因此沉淀中的溶酶体
大多是完整的。
\textbf{24.} 水为低渗液：\omterm{def:b1:cell-unit-of-life:prokeuk}{细胞器}膨胀并破裂。
线粒体和溶酶体破裂，因此沉淀 2 缩小；
琥珀酸脱氢酶（膜结合酶）部分以膜碎片的形式进入沉淀 3；
酸性磷酸酶（溶酶体内的可溶性酶）被释放
进入最终的上清液。
\textbf{25.} 纯化倍数为 $6$，产率为 \qty{60}{\%}。该倍数
受污染限制（溶酶体、过氧化物酶体、截留的\omterm{def:b1:cell-unit-of-life:cell}{细胞}溶胶：
沉淀 2 中线粒体占 \qty{60}{\%}）；产率则受线粒体损失限制，
这些线粒体或留在沉淀 1 中（被截留），或在\omterm{def:b1:cell-unit-of-life:fractionation}{匀浆}时被破坏。
\end{solution}
